\SourcePage{126}{129}
\section{Parity of a state}

Besides translations and rotations of the coordinate system (invariance under
which expresses, respectively, the homogeneity and isotropy of space),
there is one more operation that preserves the Hamiltonian of a closed
system: the \emph{inversion transformation}, which reverses the sign of
every coordinate at once (equivalently, it flips the direction of all
axes, turning a right-handed frame into a left-handed one and
conversely). Invariance of the Hamiltonian under inversion is the
mathematical expression of the mirror symmetry of space\footnote{Inversion
likewise leaves invariant the Hamiltonian of particles moving in a centrally
symmetric field (provided the coordinate origin coincides with the center of
that field).}. In classical
mechanics, the invariance of the Hamilton function under inversion gives no
\SourcePage{127}{130}
new conservation laws. In quantum mechanics, however, the situation is
essentially different.

Introduce the inversion operator $\hat P$\footnote{From the English word
\emph{parity}.}, whose action on a wave function $\psi(\mathbf r)$ consists
in reversing the signs of the coordinates:
\begin{equation}
  \hat P\psi(\mathbf r)=\psi(-\mathbf r).
  \tag{30.1}
\end{equation}
The eigenvalues $P$ of this operator, defined by the equation
\begin{equation}
  \hat P\psi(\mathbf r)=P\psi(\mathbf r),
  \tag{30.2}
\end{equation}
are easily found. Two successive applications of the inversion operator give
the identity: the arguments of the function are left entirely unchanged. In
other words, $\hat P^2\psi=P^2\psi=\psi$, so that $P^2=1$, whence
\begin{equation}
  P=\pm1.
  \tag{30.3}
\end{equation}

An eigenfunction of the inversion operator is therefore either left entirely
unaltered or acquires an overall sign change. Wave functions (and the
corresponding states) of the first type are called \emph{even}, and those
of the second type \emph{odd}.

The invariance of the Hamiltonian under inversion (that is, the commutation of
$\hat H$ and $\hat P$) therefore expresses the \emph{law of conservation of
parity}: if a state of a closed system has definite parity (if it is even or
odd), that parity is conserved in time\footnote{To avoid misunderstanding, we
recall that the discussion concerns nonrelativistic theory. Interactions exist
in nature, and are treated in relativistic theory, that violate parity
conservation.}.

The angular-momentum operator is likewise invariant under inversion:
inversion reverses the signs both of the coordinates and of the operators for
differentiation with respect to them, and hence the operator (26.2) remains
unchanged. Put differently, inversion and angular momentum commute, so that
a definite parity can coexist with definite values of $L$ and of the
projection $M$. Moreover, all states that differ only in $M$ have the same
parity. This is already evident from the independence of the properties of a
closed system from its orientation in space; formally, it can be proved from
the commutation relation $\hat L_+\hat P-\hat P\hat L_+=0$ by the same method
used to obtain (29.3) from (29.2).
\SourcePage{128}{131}
There are definite \emph{parity selection rules} for the matrix elements of
various physical quantities.

Consider scalar quantities first. Here one must distinguish \emph{true
scalars}, which remain entirely unchanged under inversion, from
\emph{pseudoscalars}, which change sign under inversion (the scalar product of
an axial vector and a polar vector is a pseudoscalar). The operator $\hat f$
of a true scalar commutes with $\hat P$. It follows that if the matrix of $P$
is diagonal, then $f$ is likewise diagonal with respect to parity: matrix
elements vanish except for the $g\to g$ and $u\to u$ transitions (here $g$
and $u$ label even and odd states).
For the operator of a pseudoscalar quantity, on the other hand,
$\hat P\hat f=-\hat f\hat P$; the operators $\hat P$ and $\hat f$
“anticommute.” The matrix element of this equality for a transition
$g\to g$ is
\[
  P_{gg}f_{gg}=-f_{gg}P_{gg},
\]
and since $P_{gg}=1$, we have $f_{gg}=0$; in the same way one finds
$f_{uu}=0$. Thus the matrix of a pseudoscalar quantity has nonzero elements
only for transitions in which the parity changes. The selection rules for
matrix elements of scalar quantities are therefore
\begin{equation}
 \begin{aligned}
  \text{true scalars:}&\qquad g\to g,\quad u\to u,\\
  \text{pseudoscalars:}&\qquad g\to u,\quad u\to g.
 \end{aligned}
 \tag{30.4}
\end{equation}

These rules can also be obtained in another way, directly from the definition
of matrix elements. Consider, for example, the integral
$f_{ug}=\int\psi_u^*\hat f\psi_g\,dq$, where $\psi_g$ is even and $\psi_u$
is odd. On reversing the signs of all coordinates, the integrand changes sign
if $f$ is a true scalar. On the other hand, an integral over all space cannot
change when the variables of integration are merely renamed. Hence
$f_{ug}=-f_{ug}$, so that $f_{ug}=0$.

The selection rules for vector quantities are obtained in the same way. One
must remember that ordinary, polar vectors reverse sign under inversion,
whereas axial vectors remain unchanged (an example is the angular-momentum
vector, the cross product $\mathbf r\times\mathbf p$ of two polar vectors).
Taking this into account, we find
\begin{equation}
 \begin{aligned}
  \text{polar vectors:}&\qquad g\to u,\quad u\to g,\\
  \text{axial vectors:}&\qquad g\to g,\quad u\to u.
 \end{aligned}
 \tag{30.5}
\end{equation}
Let us determine the parity of a one-particle state with angular momentum
$l$. Inversion ($x\to-x$, $y\to-y$, $z\to-z$) consists,
\SourcePage{129}{132}
in spherical coordinates, of the transformation
\begin{equation}
  r\to r,\qquad \theta\to\pi-\theta,\qquad
  \varphi\to\varphi+\pi.
  \tag{30.6}
\end{equation}
The angular dependence of the particle's wave function is given by the
spherical function $Y_{lm}$, which, apart from a constant immaterial here,
has the form $P_l^m(\cos\theta)e^{im\varphi}$. When $\varphi$ is replaced by
$\varphi+\pi$, the factor $e^{im\varphi}$ is multiplied by $(-1)^m$; when
$\theta$ is replaced by $\pi-\theta$, $P_l^m(\cos\theta)$ becomes
$P_l^m(-\cos\theta)=(-1)^{l-m}P_l^m(\cos\theta)$. Thus the entire function is
multiplied by $(-1)^l$ (which is independent of $m$, in agreement with the
statement above); hence the parity of a state with a given $l$ is
\begin{equation}
  P=(-1)^l.
  \tag{30.7}
\end{equation}
Accordingly, even values of $l$ yield even states and odd values yield odd
ones.

For a vector physical quantity associated with a single particle, nonzero
matrix elements can arise only in transitions with $l\to l,l\pm1$ (\secref{29}). Combining
this fact with (30.7) and with the preceding statements concerning parity
changes in vector matrix elements, we conclude that matrix elements of vector
quantities belonging to an individual particle are nonzero only for the
transitions
\begin{equation}
 \begin{aligned}
  \text{polar vectors:}&\qquad l\to l\pm1,\\
  \text{axial vectors:}&\qquad l\to l.
 \end{aligned}
 \tag{30.8}
\end{equation}
